\chapter{ABSTRACT}

\indent\indent In the era of digital transformation, the automatic classification of Intangible Cultural Heritage (ICH) data is increasingly important for cultural preservation and promotion. However, traditional CNN-based methods that depend mainly on visual information still face many challenges, especially because ICH data are multimodal and often visually similar in terms of costumes, rituals, and festival settings. To overcome these issues, I propose a deep learning approach that combines a Multimodal Graph Attention Network (MGAT) with the FastKAN model. This architecture allows images, text, and audio to be processed together in a shared feature space. Experiments on a self-collected ICH dataset show that the proposed method performs noticeably better than the baseline models. The best results are achieved by the DenseNet121 + FastKAN combination, reaching an accuracy of 94.35\%. Several smaller models, such as MobileNetV3 Small and ShuffleNetV2 1.0x, also gain improvements of more than 8\%. The results further indicate that FastKAN helps the model capture complex non-linear patterns more effectively than traditional MLPs. The practicality of the work is demonstrated through the development of the “ICH Predictor” web application, which supports easy access and visualization of the ICH classification process.\vspace{0.3cm}

\textbf{Keywords:} Intangible Cultural Heritage, Multimodal Classification, Graph Neural Networks, MGAT, FastKAN, Deep Learning.